\section{Normal cylinders and signed excess}\label{sec:normal-inputs}

The next construction compares the varifold with a smooth reference disk,
not necessarily a plane. We fix the geometric conventions and state the
two measure estimates from \cite{AE} used throughout the proof. The local
minimal-fitting statements will be introduced when the reference disk is
constructed in Section~\ref{sec:fitting}.

\subsection{Geometry and units}
Let $S$ be a small-slope $C^J$ graph over an $m$-plane, with a fixed normal
tube and uniform bounds through a sufficiently large finite order $J$.
In that tube write
\[
 X=y+z,\qquad y=\pi(X)\in S,\qquad z\in N_yS.
\]
The horizontal and vertical subspaces are $T_yS$ and $N_yS$. Let $P_T$
denote orthogonal projection onto the approximate tangent plane of the
varifold at $X$. For $v\in N_yS$, the shape operator $\mathcal S_v$ is
defined by
\[
 \langle\mathcal S_v a,b\rangle
   =\langle\II_S(a,b),v\rangle,
 \qquad a,b\in T_yS.
\]
Set
\begin{equation}\label{eq:normal-tensors}
 B=I-(P_T)_{hh},\qquad e=\tr B=\tr(P_T)_{vv},
 \qquad J_z=I-\mathcal S_z,\qquad K_z=J_z^{-1}.
\end{equation}
The tube is chosen so that $J_z$ and $K_z$ are uniformly bounded.
The tensor $B$ is positive semidefinite, $0\le e\le m$, and
\begin{equation}\label{eq:projection-blocks}
 (P_T)_{hv}(P_T)_{vh}=B-B^2,
 \qquad |(P_T)_{vh}a|^2\le B[a,a],
 \qquad |(P_T)_{vh}|\le C\sqrt e.
\end{equation}
These follow by taking blocks in $P_T^2=P_T$.

We restrict to an interior cylinder whose support is separated from its
artificial vertical boundary above every compact subset of the base.
Every field has compact base support and is multiplied by a vertical
cutoff equal to one near the relevant support. Thus the restriction has
no vertical boundary term in its first variation. All enlargements below
stay inside this cylinder.

For the positive integer $Q$ being considered, define measures on $S$ by
\begin{equation}\label{eq:normal-measures}
 M=\pi_\#\mu-Q\vol_S,\qquad D=\pi_\#(e\mu).
\end{equation}
Here $M$ is signed and $D$ is nonnegative. In a specified physical system
of Euclidean units, let $z_{\mathrm{phys}}$ denote the unscaled normal
displacement and put
\begin{equation}\label{eq:physical-error}
 \Lambda=\pi_\#(\lambda\mu),\qquad
 \lambda=|z_{\mathrm{phys}}|^2+|\mathcal H_S|^2
                       +|\nabla^\perp\mathcal H_S|^2,
 \qquad \mathcal H_S=\tr\II_S.
\end{equation}
The curvature and its derivative in this formula are evaluated in those
physical units. On scaling a base chart of physical radius $l$ to unit
radius, all measures, including reference volume, are pushed forward and
divided by $l^m$. The factor $\lambda$ is still evaluated physically; it
is not replaced by the squared rescaled height or the rescaled curvature.
For a flat reference plane the estimates of \cite{AE} use
$\Lambda=0$. Charging an additional nonnegative height term in an upper
bound does not weaken their applicability.

Normal frames are transported radially from the chart center, and base
coordinates are normalized to have Euclidean metric at that center.
For an interior subchart using the inherited frames, the small geometric
parameter remains the original
\begin{equation}\label{eq:physical-gamma}
 \gamma=l+\sup|z_{\mathrm{phys}}|.
\end{equation}
This convention matters when the signed-excess theorem is used on much
smaller cells: changing the coordinates alone does not improve a
physical curvature error.

\subsection{The two imported measure estimates}
The first statement identifies an integer using a smooth projected mass
average, then controls the total variation of $M$. The projected count
$N(y)$ counts normal-fiber entries with their integer multiplicities at
almost every base point; exceptional projection fibers are immaterial
for the statement.

\begin{lemma}[Scalar mass and propagation of the integer]\label{lem:scalar-input}
Fix dimensions, a bound for $Q$, bounded chart geometry, a mass bound, and
an interior enlargement factor $c>1$. Suppose
\[
 s^{-m}\pi_\#\mu(B_{cs})\le C_0,
 \qquad s^{-m}(D+\Lambda)(B_{cs})\le\epsilon
\]
with $\epsilon$ sufficiently small. Put $b=(1+c)/2$. On the
intermediate ball $B_{bs}$, the count selects an integer $q\ge0$ with
$\vol_S\{y\in B_{bs}:N(y)\ne q\}\le C\epsilon s^m$.
There is a fixed finite test order $J_{\mathrm{sc}}$ such that, for
every smooth scalar test $\eta$ supported in $B_{bs}$,
\begin{equation}\label{eq:scalar-distribution-before-count}
 \begin{aligned}
 \left|\int\eta\,\dd(\pi_\#\mu-q\vol_S)\right|
 &\le C\epsilon s^m\|\eta\|_{C^{J_{\mathrm{sc}}}_s},\\
 \|\eta\|_{C^{J_{\mathrm{sc}}}_s}
 &=\max_{0\le j\le J_{\mathrm{sc}}}s^j\|D^j\eta\|_\infty.
 \end{aligned}
\end{equation}
The derivatives are in base coordinates; this test order is included
among the finite orders fixed for the argument.
For every nonnegative smooth test $\chi$ supported in $B_{bs}$,
with reference integral one and
$s^m\|\chi\|_{C^{J_{\mathrm{sc}}}_s}\le C_1$ for a fixed constant
$C_1$, one also has
\begin{equation}\label{eq:scalar-average-before-count}
 \left|\int\chi\,\dd(\pi_\#\mu)-q\right|\le C\epsilon.
\end{equation}
Here the constant may depend on $C_1$; the factor $s^m$ accounts for
the size of a test whose integral, rather than its supremum, is one.
In particular, if such an average is within $1/4$ of $Q$, then $q=Q$,
and
\begin{equation}\label{eq:scalar-input}
 |M|(B_s)+\vol_S\{y\in B_s:N(y)\ne Q\}
                     \le C(D+\Lambda)(B_{cs}).
\end{equation}
The same integer propagates along comparable working balls $B_s$
satisfying these bounds, provided their overlaps have fixed positive
relative volume. The smallness threshold depends on that overlap fraction.
\end{lemma}

This is \cite[Lemma~2.5 and its proof]{AE}, including its distributional,
smooth-average, and overlap conclusions. The distributional estimate
first compares the projected measure with a scalar constant; the
integer-count argument places that constant within $C\epsilon$ of $q$.
It therefore gives Equation~\eqref{eq:scalar-distribution-before-count}
before any identification with $Q$. The enlargement $c$ can be any
fixed number greater than one, by a finite interior covering. In
particular we use $c=2$ when the room has already been fixed. The integer
must be identified before applying Equation~\eqref{eq:scalar-input}; a
bound for the tilt alone does not specify which integer to subtract.

For a positive integer $j$, write $\exp_j$ and $\log_j$ for $j$-fold
composition. For sufficiently large $T$, let
\begin{equation}\label{eq:signed-rate}
 F_Q(T)=\exp_{2Q}\bigl(\sqrt{\log_{2Q}T}\bigr).
\end{equation}
We will use only the properties
\[
 (\log T)^a=o(F_Q(T))\quad\text{for every fixed }a>0,
 \qquad F_Q(T)=T^{o(1)}.
\]

\begin{theorem}[Signed excess]\label{thm:signed-input}
Fix an integer $Q\ge1$ and $\kappa>0$. In the normal-cylinder setting above, scale a
physical base radius $l$ to one and suppose $B_4$ is available. Set
$\tau=|M|+D+\Lambda$. If
\begin{equation}\label{eq:signed-hypotheses}
 \tau(B_4)\le d=e^{-T},\qquad
 l+\sup|z_{\mathrm{phys}}|\le d^\kappa,
\end{equation}
where the supremum is over the restricted support above the available
chart, then, for $T$ sufficiently large,
\begin{equation}\label{eq:signed-input}
 \int w\,\dd(2M-D-C_Q\Lambda)
       \le d\exp[-F_Q(T)]
\end{equation}
whenever $w$ is supported in $B_1$, $0\le w\le1$, and $\Lip w\le1$.
The finite derivative order $J$, threshold, and constants depend only
on the stated fixed data. For a subchart in inherited frames, the second
condition is imposed on $\gamma$ in Equation~\eqref{eq:physical-gamma}.
For a flat reference, take $\Lambda=0$ and omit the second condition.
\end{theorem}

This is \cite[Theorem~3.1]{AE}. Multiplying the test by a fixed constant
allows any fixed Lipschitz bound, with that constant on the right-hand
side. The theorem does not assume that the multiplicity is bounded by
$Q$: $Q$ is the integer subtracted in $M$.

\subsection{Fixed parameters}
All dimensional constants and interior enlargement factors are fixed
before small scales are sent to zero. Fix an integer kernel exponent $p\ge4$, put
$J_0=m+p+4$, and choose
\begin{equation}\label{eq:parameter-order}
 0<\beta_1<\tfrac14,\qquad
 0<\zeta<\min\{1,1/m\},\qquad \kappa>0,\qquad
 \kappa(2J_0+4)<2\beta_1,\qquad 3\kappa<\zeta.
\end{equation}
We then fix derivative orders large enough for Theorem~\ref{thm:signed-input}
and the compactness arguments below. Finally we choose the fitting order
and its small height coefficient in Section~\ref{sec:fitting}.
This order is noncircular: $J_0$, $\beta_1$, $\zeta$, and $\kappa$ do not
depend on the fitting order. The two inequalities in
Equation~\eqref{eq:parameter-order} will respectively permit the signed
comparison in the frequency estimate and in the height-normalized limit.
